\documentclass[10pt,a4paper]{amsart}
\usepackage[margin=19mm]{geometry}
\usepackage{amsmath,amssymb,mathtools}
\usepackage[scr=rsfs,cal=euler]{mathalfa}
\usepackage{xfrac}
\usepackage[unicode,hidelinks,bookmarks=false]{hyperref}
\newcommand{\ie}{i.e.\ }
\newcommand{\cf}{cf.\ }
\newcommand{\resp}{respectively\ }
\newcommand{\Subsection}{\subsection}
\providecommand{\nobreakdash}{\nobreak\hbox{-}}
\setlength{\emergencystretch}{2em}
\multlinegap0pt
\usepackage{amssymb}
\usepackage[shortlabels]{enumitem}
\usepackage{mathtools}
\usepackage{tikz}
\newtheorem{theorem}{Theorem}
\newtheorem{lemma}{Lemma}
\newcommand{\Om}{\Omega}
\newcommand{\dO}{\delta\Omega}
\title{Effect of edge-stretching on Steklov eigenvalues \\and sharp Steklov eigenvalue bounds on leaf--boundary trees}
\author{Jiangdong Ai, Yizhe Ji, Xiaopan Lian, Kun Yang}
\date{Source: arXiv:2603.26251v1 (CC BY 4.0)}
\begin{document}
\maketitle

\noindent\emph{Source and licence.} Effect of edge-stretching on Steklov eigenvalues \\and sharp Steklov eigenvalue bounds on leaf--boundary trees. Version arXiv:2603.26251v1 (CC BY 4.0), distributed by arXiv under the Creative Commons Attribution 4.0 International licence, http://creativecommons.org/licenses/by/4.0/, as recorded in the arXiv licence metadata for this exact version and shown on the abstract page https://arxiv.org/abs/2603.26251v1. Full licence notice: \texttt{licenses/arxiv-2603.26251-CC-BY-4.0.txt}.\par\smallskip

\noindent\textit{Editorial scope.} Counted unit: the article's theorem thm:lambda2sharp (source0.tex lines 234--242). The article's complete original proof of it is delivered with it (source0.tex lines 616--659, one \texttt{proof} environment of 44 lines, which the article places away from the statement). No mathematics is written, repaired or re-derived: the statement and the proof are the article's own lines, in the article's own notation, and no retained line is substituted or re-flowed.  Retained uncounted as the source-local context the proof needs: lines 573--576; lines 583--592.  Nothing was omitted from any retained span, so the package carries no omission record.  No retained line contains a citation, so the wrapper carries no bibliography and no bibliography entry.  No retained line contains a figure environment, an image-inclusion command or drawing code, so no image is delivered and the package declares no asset dependency.  Editorial formatting only, in addition: an AMS \texttt{amsart} preamble that restates the article's own declarations and macros used by the retained lines, namely the environments \texttt{theorem}, \texttt{lemma} on separate counters, together with the standard packages those lines need.  The retained environments print in this document as Theorem 1, Lemma 1 in order of appearance, whereas the article prints the same items with the numbering of its own full document.  The complete native source of the article as distributed in the arXiv e-print for this exact version is bundled unchanged as \texttt{sources/197268/source0.tex}.
\begin{lemma}\label{lem:centroid}
Let $T$ be a tree and $\dO$ its set of leaves.
There exists a vertex $v$ such that each component $C$ of $T\setminus\{v\}$ satisfies $|\dO\cap C|\le |\dO|/2$.
\end{lemma}

\begin{lemma}\label{lem:vertex2value}
Let $T$ be a tree with leaf boundary $\dO$ and $|\dO|=\ell$.
Fix a vertex $v$ of degree $r\ge 2$ and let $C_1,\dots,C_r$ be the components of $T\setminus\{v\}$.
Write $b_i:=|\dO\cap C_i|$ and fix an index $i_0$.
Set $s:=b_{i_0}$ and assume $0<s\le\ell/2$.
Then
\begin{equation}\label{eq:vertex2value}
\lambda_2(T,\dO)\le \frac{r-1}{r}\cdot\frac{\ell}{s(\ell-s)}.
\end{equation}
\end{lemma}

\begin{theorem}\label{thm:lambda2sharp}
Let $T=(V,E)$ be a tree with leaf boundary $\dO$.
Let $D\ge 2$ be the maximum degree of $T$ and let $\ell=|\dO|$ be the number of leaves.
Then
\begin{equation}\label{eq:mainbound_intro}
\lambda_2\le \frac{D}{\ell}.
\end{equation}
Moreover, equality holds if and only if $T$ is a star.
\end{theorem}

\begin{proof}[Proof of Theorem~\ref{thm:lambda2sharp}]
Let $v$ be a leaf centroid as in Lemma~\ref{lem:centroid}.
Let $C_1,\dots,C_r$ be the components of $T\setminus\{v\}$, so $r=\deg(v)\le D$.
Set $b_i:=|\dO\cap C_i|$.
By the centroid property, each $b_i\le\ell/2$.
Let $i_0$ maximize $b_i$ and write $s:=b_{i_0}$.
By pigeonhole, $s\ge\ell/r$.

Apply Lemma~\ref{lem:vertex2value} at $v$ and $C_{i_0}$ to obtain
$$
\lambda_2\le \frac{r-1}{r}\cdot\frac{\ell}{s(\ell-s)}.
$$
Since $s\le\ell/2$, the function $x\mapsto\ell/(x(\ell-x))$ is decreasing on $(0,\ell/2]$, hence
$$
\frac{\ell}{s(\ell-s)}\le \frac{\ell}{(\ell/r)(\ell-\ell/r)}=\frac{r^2}{\ell(r-1)}.
$$
Therefore $\lambda_2\le \frac{r-1}{r}\cdot\frac{r^2}{\ell(r-1)}=\frac{r}{\ell}\le\frac{D}{\ell}$.

For the equality statement, equality throughout forces
$
r=D, s=\frac{\ell}{r},
$
and the harmonic extension must have the same energy as the piecewise--constant extension used in
Lemma~\ref{lem:vertex2value}. Hence the function $f$ constructed in Lemma~\ref{lem:vertex2value}, which is constant on each component
$C_i$, must itself be harmonic on $\Om$.

We claim that each component $C_i$ contains no interior vertex. Indeed, suppose that some $C_i$
contains an interior vertex. Since $C_i$ is a component of $T\setminus\{v\}$, there is a unique vertex
$w\in C_i$ adjacent to $v$. This vertex $w$ lies in $\Om$, because all boundary vertices are leaves
and $w$ has the neighbor $v$ together with at least one neighbor inside $C_i$. By construction, $f$
is constant on $C_i$, so every neighbor of $w$ inside $C_i$ has the same value as $f(w)$. On the other
hand, the value at $v$ is different from the constant value on $C_i$. Therefore
$$
\sum_{w'\sim w}\bigl(f(w)-f(w')\bigr)=f(w)-f(v)\neq 0,
$$
so $f$ is not harmonic at $w$, a contradiction.

Thus every $C_i$ contains no interior vertex, and hence every $C_i$ consists of a single boundary
leaf. Therefore $T$ is a star. Conversely, for the star $K_{1,\ell}$ we have
$
\lambda_2=1=\frac{D}{\ell},
$
since $D=\ell$. Hence equality holds in~\eqref{eq:mainbound_intro}.
\end{proof}

\end{document}
